\documentclass[10pt,aspectratio=169]{beamer}

% --- Minimal theme ---------------------------------------------------------
\usetheme{default}
\usecolortheme{dove}
\usefonttheme{professionalfonts}
\setbeamertemplate{navigation symbols}{}
\setbeamertemplate{footline}[frame number]
\setbeamertemplate{itemize items}[circle]
\setbeamertemplate{section in toc}[sections numbered]
\setbeamerfont{frametitle}{size=\large,series=\bfseries}
\setbeamercolor{frametitle}{fg=black}
\setbeamercolor{title}{fg=black}
\setbeamercolor{structure}{fg=black}

\usepackage{booktabs}
\usepackage{array}
\usepackage{graphicx}
\graphicspath{{../figs/}}

\title{Which reanalysis is ``best''?}
\subtitle{An out-of-sample predictive test of five temperature products}
\author{}
\date{\today}

\begin{document}

\begin{frame}
  \titlepage
\end{frame}

\begin{frame}{Outline}
  \tableofcontents
\end{frame}

% ===========================================================================
\section{Motivation}

\begin{frame}{The previous deck left us stuck}
  \begin{itemize}
    \item Five reanalyses, one panel, same specification $\Rightarrow$ five
          different ME curves.
    \item Pairwise L1 distances between MEs range from 0.0005 (ERA5 vs JRA)
          to 0.0035 (ERA5 vs MERRA2).
    \item The fits themselves cannot tell us which is ``right'' --- each
          regression is internally consistent.
  \end{itemize}
  \vspace{0.5em}
  \textbf{Question:} does any dataset's ME predict growth better than the others?
\end{frame}

% ===========================================================================
\section{Method}

\begin{frame}{Frisch--Waugh--Lovell + k-fold on the residuals}
  For each dataset $k$, fit the nuisance model on the full (balanced) sample:
  \[
    y = \boldsymbol\gamma^\top \mathbf{p} + \alpha_i + \tau_t
      + \delta_i t + \phi_i t^2 + \varepsilon \qquad(\mathbf{p} = \text{precip.})
  \]
  Repeat separately for $y = \texttt{growthWDI}$, $y = \texttt{temp\_1}$,
  $y = \texttt{temp\_1}^2$ and keep the residuals $\tilde y, \tilde T, \tilde{T^2}$.

  \vspace{0.4em}
  On these residuals, the temperature coefficients reproduce the original
  regression (FWL). The \emph{out-of-sample} step is then a small linear fit:
  \[
    \tilde y \;=\; \beta_1 \tilde T + \beta_2 \tilde{T^2} + u
  \]
\end{frame}

\begin{frame}{Three cross-validation schemes}
  \begin{description}
    \item[Random 10-fold] --- partition on (iso, year). Pure predictive fit.
    \item[Leave-country-out] --- hold out every country in turn. Tests whether
          the ME extrapolates to \emph{new countries} (158 folds).
    \item[Leave-year-out] --- hold out every year in turn. Tests temporal
          transfer (31 folds, 1980--2013 after balancing).
  \end{description}
  \vspace{0.3em}
  For each fold we use \texttt{.lm.fit} (Base R, QR only) --- benchmarked at
  $\sim$140\,\textmu s, versus $\sim$550\,\textmu s for \texttt{lm} and
  $\sim$2.8\,ms for \texttt{feols}. A few thousand fits run in under a second.
\end{frame}

\begin{frame}{Scoring}
  \begin{itemize}
    \item Per fold: \textbf{MSE} of $\tilde y - \hat{\tilde y}$ on held-out rows.
    \item Per (dataset, split): mean MSE, standard error across folds, and
          \[
            R^2_{\text{oos}} \;=\; 1 - \frac{\overline{\mathrm{MSE}}}{\widehat{\mathrm{Var}}(\tilde y)}.
          \]
    \item Same random seed across datasets $\Rightarrow$ identical 10-fold
          partition $\Rightarrow$ comparisons are paired at the fold level.
  \end{itemize}
\end{frame}

% ===========================================================================
\section{Results}

\begin{frame}{Headline: datasets are indistinguishable}
  \centering
  \includegraphics[width=\linewidth]{oos_mse.pdf}

  \vspace{0.3em}
  \footnotesize
  Error bars = $\pm 1$ SE across folds. Scale of between-dataset differences
  is $\sim 10^{-5}$; within-dataset fold variation is $10\times$ larger.
\end{frame}

\begin{frame}{Same picture in OOS $R^2$}
  \centering
  \includegraphics[width=\linewidth]{oos_r2.pdf}

  \vspace{0.3em}
  \footnotesize
  Under random 10-fold all five datasets give $R^2_{\text{oos}} \in
  [0.003, 0.007]$. Under leave-country-out \emph{all five} are negative
  ($\approx -0.06$): the ME does not generalise to unseen countries.
\end{frame}

\begin{frame}{Numbers}
  \centering
  \small
  \begin{tabular}{l l r r r r}
    \toprule
    Dataset & Split & CV-MSE & SE & $R^2_\text{oos}$ & Folds \\
    \midrule
    GMFD      & random-10fold     & 0.00269 & 0.00018 & \textbf{0.00732} & 10  \\
    JRA-3Q    & random-10fold     & 0.00270 & 0.00018 & 0.00542          & 10  \\
    ERA5-025  & random-10fold     & 0.00271 & 0.00018 & 0.00391          & 10  \\
    UDel      & random-10fold     & 0.00271 & 0.00018 & 0.00392          & 10  \\
    MERRA2    & random-10fold     & 0.00271 & 0.00018 & 0.00290          & 10  \\
    \midrule
    GMFD      & leave-country-out & 0.00288 & 0.00046 & $-0.0608$        & 158 \\
    JRA-3Q    & leave-country-out & 0.00289 & 0.00046 & $-0.0634$        & 158 \\
    ERA5-025  & leave-country-out & 0.00289 & 0.00046 & $-0.0644$        & 158 \\
    UDel      & leave-country-out & 0.00289 & 0.00046 & $-0.0645$        & 158 \\
    MERRA2    & leave-country-out & 0.00289 & 0.00046 & $-0.0655$        & 158 \\
    \midrule
    GMFD      & leave-year-out    & 0.00270 & 0.00033 & 0.00575          & 31  \\
    JRA-3Q    & leave-year-out    & 0.00270 & 0.00033 & 0.00479          & 31  \\
    UDel      & leave-year-out    & 0.00271 & 0.00033 & 0.00334          & 31  \\
    ERA5-025  & leave-year-out    & 0.00271 & 0.00033 & 0.00326          & 31  \\
    MERRA2    & leave-year-out    & 0.00271 & 0.00033 & 0.00240          & 31  \\
    \bottomrule
  \end{tabular}
\end{frame}

\begin{frame}{Consistent ranking (but tiny)}
  Across all three splits the same order emerges:
  \[
    \text{GMFD} \;\succeq\; \text{JRA-3Q} \;\succeq\; \text{ERA5-025, UDel}
                    \;\succeq\; \text{MERRA2}.
  \]
  \vspace{0.3em}
  But the MSE gap from best (GMFD) to worst (MERRA2) is
  $\approx 2 \times 10^{-5}$, vs.\ per-fold SE $\sim 2 \times 10^{-4}$. The
  ranking is \emph{suggestive}, not statistically separable.
\end{frame}

% ===========================================================================
\section{Discussion}

\begin{frame}{What this tells us}
  \begin{itemize}
    \item For predicting held-out growth, the choice of reanalysis is
          \textbf{practically irrelevant}. Different MEs, identical predictive fit.
    \item The ME \emph{fails to generalise across countries}
          ($R^2_{\text{oos}} \approx -0.06$ under leave-country-out).
          Any apparent predictive power at the (iso, year) level is coming
          from within-country variation, not from a stable cross-country
          temperature--growth relationship.
    \item Reanalysis choice matters for the \emph{shape} of the damage function
          (our earlier finding), but not for how well the function
          \emph{describes} growth in the sample.
  \end{itemize}
\end{frame}

\begin{frame}{Caveats}
  \begin{itemize}
    \item \textbf{Residualisation leakage}: FEs see test data during the
          partial-out step. Symmetric across datasets, so the horse race is
          fair; but absolute $R^2_{\text{oos}}$ values are slightly optimistic.
          Re-doing residualisation inside each fold is a natural follow-up.
    \item \textbf{Low power}: country-year growth is very noisy.
          Differences of $10^{-5}$ in MSE may exist but are below our
          resolution.
    \item \textbf{Balanced sample only}: we restrict to 4{,}592 (iso, year)
          pairs observed in all five datasets (1980--2013).
  \end{itemize}
\end{frame}

\begin{frame}{Takeaway}
  \begin{center}\Large
    We cannot pick a ``best'' reanalysis from predictive fit alone.\\[0.6em]
    Any selection must appeal to \emph{external} criteria:\\[0.3em]
    station-benchmark agreement, physical plausibility, coverage.
  \end{center}
\end{frame}

\end{document}
